\section{Introduction}
A finite observer does not receive a geometric object to be compressed. It executes an experiment: it chooses an action, pays for a report, retains a label, and eventually commits to a prediction before a protected test is performed. The distribution to be approximated is created by those choices. Even when the current test response is exactly visible, a stopping rule changes its law, and the cost of obtaining a favorable response cannot be removed by conditioning.

We study the passage
\[
\begin{aligned}
&\text{prepared kernels}\longrightarrow\text{executable tests}
 \longrightarrow\text{predictive realization}\\
&\hspace{2em}\longrightarrow\text{actual acquisition and terminal mass}
 \longrightarrow\text{finite-resource risk}.
\end{aligned}
\]
The principal theorem identifies a class on which this passage is exact. Each acquisition action has a prescribed current-response law, but different actions need not have the same law, support, cost, or local dimension. We call this property \emph{task regeneration}; it is a condition on a specified executable test, not an assertion that all hidden physical variables are independent or reset. Theorem~\ref{thm:central} converts every history-dependent acquisition and stopping procedure into a scalar action-occupation profile, and conversely implements every feasible profile with one live state. Optimizing a codebook and optimizing the acquired law are therefore one constrained problem. The full causal memory price is exactly one additional live label.

Two issues make the result different from static quantization. First, action counts and selected terminal submeasures must be compatible. The selected mass belonging to an action cannot exceed its expected acquisition count times its one-call law. Second, the same codebook must serve the selected mass from every action. These conditions lead to a single mass--cost dual, in which the same multiplier is tested against all actions. Optimizing each action with a separate codebook and then taking an average is not the problem considered here. A finite-dimensional extremal argument reduces the executing action mixture to two actions under one scalar budget. This support reduction is a standard linear-programming mechanism; its role here is to preserve the exact acquired law and the total observer-state count.

The exact frontier is complemented by three quantitative results. Uniform conditional domination, imposed at every surviving history before selecting a program, gives matching profile bounds for correlated controlled experiments. A raw perturbation of regenerative response laws can be repaired at the \emph{same} expected acquisition and state budgets with a linear error modulus. Finally, a mixed atomic--regular acquired law has a continuous-mass factor that vanishes at an explicit acquisition threshold. Its rate changes with the number of output labels, even though the physical carrier is fixed.

Theorem~\ref{thm:gaussian} supplies an informative-filter consequence outside task regeneration. For a symmetric positive two-state hidden chain with Gaussian reports, it characterizes the optimum among \emph{all} Borel randomized two-state autonomous observers. The proof combines the actual pre-report occupation equations with the classical likelihood-ratio ordering. A one-parameter threshold determines both the exact mean acquisition and the attained prediction risk. A separately certified construction compares this optimum with checkpoint procedures for the same physical stopping experiment; it is not described as the unknown global checkpoint optimum.

The foundations and realizations remain distinct. Theorem~\ref{thm:central} is proved before the classical and quantum kernels are introduced. Those kernels verify its conditional-domination hypotheses directly. The Gaussian theorem uses the more general occupation representation, not a false assertion that a Gaussian posterior is recovered from its last report. Full proofs of the general common-program occupation criterion, the informative Bayesian three-defect identity, global-cover stability bounds, and earlier singular and finite-budget results are retained in the technical companion~\cite{GTF49,GTF32}. The present theorem does not replace them by a narrower definition.

\subsection*{Relation to earlier mathematics}
Population trimming and optimal quantization already optimize a density bounded by a reference measure~\cite{CGM}. Budgeted rejection sampling optimizes acceptance under a sampling constraint~\cite{Verine}; recent finite-sample selection results address a different fixed-sample architecture~\cite{UR}. Stationary policy reductions and finite mixtures under occupation constraints belong to the established theory of constrained Markov decision processes~\cite{DPR,DPRchat}. We use these mechanisms with explicit attribution and give direct proofs for the present setting. Our claim is not that occupation measures, rejection sampling, convex duality, or likelihood-ratio tests are new. The contributions are their exact conjunction with action-dependent acquired geometry, the one-live-state implementation of the entire attainable terminal profile, the same-budget perturbation repair, and the stated singular and informative-filter consequences.

Predictive equivalence and statistical experiment comparison have a separate classical history~\cite{SC,Fritz}. Belief compression~\cite{PB}, Bayes-risk quantization~\cite{Varshney}, and finite-memory filtering approximation~\cite{Kim} remain relevant comparisons. The present finite-state count is not a mutual-information rate, a number of stored floating-point coordinates, or a bound on temporary computation. No claim of a complete classification of causal compressibility or of a universal learning rate for unknown kernels is made.
